\section*{Bab \ref{ch:b3:remote-sensing-gis} --- Penginderaan Jauh dan SIG}

\begin{solution}{exo:b3:remote-sensing-gis:1}
Lahan: $(0.42 - 0.06)/(0.42 + 0.06) = 0.75$, vegetasi hijau rapat, tanaman
dengan daun penuh. Lahan yang dibajak: $(0.28 - 0.21)/(0.28 + 0.21) = 0.14$, hampir
gundul, sebuah tanah yang sedikit lebih cerah ke arah inframerah.
\end{solution}

\begin{solution}{exo:b3:remote-sensing-gis:2}
(a)~$185000/30 \approx 6170$ piksel; $185000/15 \approx 12300$ pada
pita pankromatik. (b)~$2^{12} = 4096$ tingkat; $2^8 = 256$. (c)~16 hari
saja; dengan Landsat~9 pada orbit yang sama selang 8 hari, 8 hari.
\end{solution}

\begin{solution}{exo:b3:remote-sensing-gis:3}
(a)~Satu pinggiran adalah $\lambda/2 = \qty{27.7}{mm}$; enam pinggiran
adalah \qty{166}{mm} perubahan sepanjang garis pandang, di sini menjauhi satelit, karena
cekungan itu mengalami penurunan. (b)~$56.6/2 = \qty{28.3}{mm}$.
\end{solution}

\begin{solution}{exo:b3:remote-sensing-gis:4}
(a)~$t = 2d/c = 2\times\num{5e5}/\num{3.00e8} = \qty{3.34}{ms}$.
(b)~$c\,\delta t/2 = \num{3.00e8}\times\num{8e-10}/2 = \qty{0.12}{m}$; permukaan
yang tersampel banyak foton jauh lebih tepat letaknya, dengan presisi
menurun sebanding dengan satu per akar jumlahnya.
\end{solution}

\begin{solution}{exo:b3:remote-sensing-gis:5}
(a)~Dalam satuan SI $L = \qty{9.0e6}{W.m^{-2}.sr^{-1}.m^{-1}}$. Dengan $c_1 =
2hc^2 = \qty{1.191e-16}{W.m^2.sr^{-1}}$, $c_2 = \qty{0.014388}{m.K}$ dan
$\lambda^5 = (\num{1.09e-5})^5 = \num{1.539e-25}$, $c_1/(\lambda^5L) = 86.0$,
dan $T_b = c_2/(\lambda\ln 87.0) = \qty{295.6}{K}$. (b)~Suhu kinetik
memenuhi $0.98\,B_\lambda(T) = L$, yakni $B_\lambda(T) = L/0.98$: rumus
yang sama memberikan $T = \qty{296.9}{K}$, lebih tinggi \qty{1.3}{K}.
\end{solution}

\begin{solution}{exo:b3:remote-sensing-gis:6}
(a)~$R = 693/\cos\ang{35} = \qty{846}{km}$. (b)~Antena nyata: $R\lambda/L =
\num{8.46e5}\times0.0555/12.3 = \qty{3.8}{km}$; apertur sintetik: $L/2 =
\qty{6.2}{m}$, sekitar 600 kali lebih halus. (c)~Apertur sintetik adalah
panjang lintasan tempat sebuah titik tetap berada dalam berkas, $R\lambda/L =
\qty{3.8}{km}$.
\end{solution}

\begin{solution}{exo:b3:remote-sensing-gis:7}
$\partial z/\partial x = (418 - 405)/60 = 0.217$ (timur dikurangi barat),
$\partial z/\partial y = (418 - 401)/60 = 0.283$ (utara dikurangi selatan);
$|\nabla z| = 0.357$, dengan kemiringan $\arctan 0.357 = \ang{19.6}$. Tanah
naik ke arah timur laut, sehingga menghadap ke barat daya: $A =
\operatorname{atan2}(-0.217, -0.283) = \ang{217}$. Bayangan lereng: $\cos\ang{45}
\cos\ang{19.6} + \sin\ang{45}\sin\ang{19.6}\cos(\ang{315} - \ang{217}) =
0.666 - 0.033 = 0.63$, sedikit lebih gelap daripada tanah datar (0.71): lereng
berbelok sedikit menjauhi sumber cahaya di barat laut.
\end{solution}

\begin{solution}{exo:b3:remote-sensing-gis:8}
Jumlah baris 52, 48, 38, 22; jumlah kolom 54, 46, 38, 22; diagonal 133 dari
160: akurasi keseluruhan \qty{83.1}{\percent}. Akurasi pengguna: air
$48/52 = 0.92$, hutan $37/48 = 0.77$, tanaman $29/38 = 0.76$, kota
$19/22 = 0.86$. Akurasi produsen: air $48/54 = 0.89$, hutan $37/46 = 0.80$,
tanaman $29/38 = 0.76$, kota $19/22 = 0.86$. $p_e = (52\cdot54 + 48\cdot46 +
38\cdot38 + 22\cdot22)/160^2 = 0.271$, $\kappa = (0.831 - 0.271)/(1 -
0.271) = 0.77$. Tanaman paling sering terlewat: 9 dari 38 piksel tanaman yang sebenarnya
masuk ke tempat lain, sebagian besar ke hutan.
\end{solution}

\begin{solution}{exo:b3:remote-sensing-gis:9}
(a)~$L = 0.04\,B_\lambda(1443) + 0.96\,B_\lambda(290)$ pada
\qty{10.9}{\micro\metre}, \qty{28.6}{W.m^{-2}.sr^{-1}.\micro\metre^{-1}},
dengan \omterm{def:b3:remote-sensing-gis:brightness}{suhu kecerahan} \qty{396}{K}, dibandingkan dengan suhu rata-rata
$0.04\times1443 + 0.96\times290 = \qty{336}{K}$. (b)~Pada
\qty{2.2}{\micro\metre}, $B_\lambda(1443)/B_\lambda(290) = \num{6.8e7}$:
lava sepenuhnya mendominasi piksel, dan \omterm{def:b3:remote-sensing-gis:brightness}{suhu kecerahannya}
\qty{845}{K}. Pita inframerah gelombang pendek mendeteksi titik panas kecil jauh
lebih baik; pita termal yang peka terhadap seluruh piksel justru
yang mengukur tanah di sekitarnya.
\end{solution}

\begin{solution}{exo:b3:remote-sensing-gis:10}
Dengan satelit di timur, $d_1 = u_E\sin\theta + u_z\cos\theta$; ke
barat, $d_2 = -u_E\sin\theta + u_z\cos\theta$. Dengan menjumlahkan dan mengurangkan,
$u_z = (d_1 + d_2)/(2\cos\theta) = 102/1.638 = \qty{62.3}{mm}$ dan $u_E =
(d_1 - d_2)/(2\sin\theta) = 26/1.147 = \qty{22.7}{mm}$. Satu lintasan memberikan satu
persamaan untuk dua besaran tak diketahui: setiap campuran pengangkatan dan gerak ke timur dengan
proyeksi yang sama pada garis pandang memberikan pinggiran yang sama.
\end{solution}

\begin{solution}{exo:b3:remote-sensing-gis:11}
(a)~Menurut \cref{thm:b3:remote-sensing-gis:fringe}, $\Delta\phi =
4\pi\Delta r/\lambda < \pi$ mensyaratkan $\Delta r < \lambda/4 =
\qty{13.9}{mm}$. (b)~$u_z = k z_0(r^2 + z_0^2)^{-3/2}$ dengan $k = (1 -
\nu)\Delta V/\pi$; $\mathrm du_z/\mathrm dr = -3kz_0r(r^2 + z_0^2)^{-5/2}$,
yang turunannya lenyap di mana $(r^2 + z_0^2) = 5r^2$, yaitu pada $r = z_0/2$. Berikut
\[
\left|\frac{\mathrm du_z}{\mathrm dr}\right| = \frac{3k}{2z_0^3}\Big(\frac54\Big)^{-5/2}
= \frac{0.859\,k}{z_0^3} = \frac{0.859\times\num{2.39e6}}{\num{6.4e10}} = \num{3.2e-5},
\]
\qty{32}{mm} per kilometer. (c)~$13.9/32 = \qty{0.43}{km}$: piksel lebih kecil
dari sekitar \qty{430}{m}, yang mudah dipenuhi Sentinel-1.
\end{solution}

\begin{solution}{exo:b3:remote-sensing-gis:12}
(a)~Untuk $N, R \ge 0$ yang tidak keduanya nol, $|N - R| \le N + R$, sehingga $|\mathrm{NDVI}|
\le 1$; mengganti $N$ dan $R$ dengan $aN$ dan $aR$ mengalikan pembilang dan
penyebut dengan $a$. (b)~Tanpa: $(0.502 - 0.051)/(0.502 + 0.051) =
0.816$; dengan: $(0.512 - 0.081)/(0.512 + 0.081) = 0.727$. (c)~Radiansi lintasan
menurunkan indeks vegetasi hijau, sekitar sepersepuluh di sini, dan
dengan besaran yang berubah sesuai kabut tipis: peta citra yang belum dikoreksi akan
menunjukkan cuaca sama banyaknya dengan vegetasi, dan perbandingan antartanggal
tidak akan bermakna.
\end{solution}

\begin{solution}{pb:b3:remote-sensing-gis:1}
\textbf{1.} Hutan $(0.38 - 0.04)/0.42 = 0.81$; aliran 1890 $0.23/0.37 =
0.62$; aliran 1960 $0.10/0.32 = 0.31$; aliran 2005 $0.02/0.26 = 0.08$.

\textbf{2.} 1890, lalu 1960, lalu 2005: aliran tertua paling hijau,
sesuai urutan umurnya; aliran itu belum mencapai hutan.

\textbf{3.} Batuan telanjang memantulkan sedikit lebih banyak di inframerah dekat daripada di
merah, seperti pasir cokelat pada gambar bab ini; lumut kerak dan tumbuhan pertama
menambah sedikit: 0.08 adalah nilai tanah yang hampir telanjang.

\textbf{4.} Lereng yang membelakangi cahaya meredupkan kedua pita dengan proporsi yang sama;
rasio menghilangkannya, sedangkan reflektansi inframerah dekat saja
akan keliru menganggap hutan yang teduh sebagai vegetasi jarang.

\textbf{5.} Berdasarkan \cref{prop:b3:remote-sensing-gis:planck-inversion} dengan
$\lambda = \qty{10.9}{\micro\metre}$: dataran \qty{291.1}{K},
piksel fumarol \qty{302.3}{K}.

\textbf{6.} $0.97\,B_\lambda(T) = \num{8.40e6}$ memberikan $T = \qty{293.0}{K}$.

\textbf{7.} $B_\lambda(\qty{370}{K}) =
\qty{22.48}{W.m^{-2}.sr^{-1}.\micro\metre^{-1}}$; $f = (9.95 - 8.40)/(22.48
- 8.40) = 0.110$, sekitar \qty{1100}{m^2} dari piksel \qty{e4}{m^2}.

\textbf{8.} Radiansi adalah fungsi cembung terhadap suhu
(\cref{prop:b3:remote-sensing-gis:mixed}): bagian yang lebih panas menambahkan radiansi lebih banyak
daripada porsi suhunya, sehingga piksel tampak lebih panas daripada rata-ratanya,
$0.110\times370 + 0.890\times291.1 = \qty{299.8}{K}$ dibandingkan dengan
\qty{302.3}{K}.

\textbf{9.} Fase setiap gema bergantung pada lintasan yang tepat; perubahan besar
sudut pandang mengubah jumlah gema dari para penghambur, dan
\omterm{def:b3:remote-sensing-gis:insar}{interferogram} kehilangan koherensinya, dan ia menambahkan fase akibat relief.
$48/12 = 4$ siklus dari satu satelit.

\textbf{10.} $5.4\times27.73 = \qty{149.7}{mm}$ menuju satelit.

\textbf{11.} $u_z = 149.7/\cos\ang{35} = \qty{182.8}{mm}$.

\textbf{12.} Tanah bergerak ke luar sekaligus ke atas; di sisi timur
gerakan ke luar mengarah ke satelit dan menambah pengangkatan, di sisi barat
gerakan itu mengurangi (\cref{prop:b3:remote-sensing-gis:mogi-los}). Perubahan itu
turun lebih cepat ke barat, tempat pinggiran interferensinya rapat; polanya
meregang ke timur.

\textbf{13.} Pada garis itu gerak horizontal adalah utara--selatan,
tegak lurus terhadap arah pandang, sehingga radar hanya melihat $u_z\cos\theta$ saja, yang
berbentuk seperti pengangkatan. Pengangkatan turun separuh pada $r = 0.766\,z_0$, sehingga $z_0 =
3.4/0.766 = \qty{4.44}{km}$.

\textbf{14.} Separuh dari 5.4: 2.7 pinggiran.

\textbf{15.} Deformasi tampak sama pada kedua \omterm{def:b3:remote-sensing-gis:insar}{interferogram} setelah
dinormalisasi dengan selang waktunya, atau terus membesar dari pasangan ke pasangan; pola
atmosfer berubah dari tanggal ke tanggal dan sering mengikuti
relief. Pasangan kedua yang menunjukkan pola lingkaran konsentris yang sama memastikan
deformasi itu.

\textbf{16.} $\Delta V = \pi u_z(0)z_0^2/(1 - \nu) = \pi\times0.1828\times
4436^2/0.75 = \qty{1.51e7}{m^3}$.

\textbf{17.} $\num{1.51e7}/48 = \qty{3.1e5}{m^3}$ per hari, \qty{1.1e8}{m^3}
per tahun jika laju itu tetap.

\textbf{18.} $u_r \propto r/(r^2 + z_0^2)^{3/2}$ terbesar di tempat
turunannya nol, $r^2 + z_0^2 = 3r^2$, $r = z_0/\sqrt2 = \qty{3.14}{km}$;
di sana $u_r = \qty{70}{mm}$, \qty{38}{\percent} dari pengangkatan pusat.

\textbf{19.} Jari-jari \qty{1.48}{km}, volume $\frac43\pi(1.48)^3 =
\qty{13.6}{km^3}$; perubahannya adalah $\num{1.51e7}/\num{1.36e10} =
\qty{0.11}{\percent}$ dari volume itu.

\textbf{20.} Tidak: sepersepuluh pinggiran adalah \qty{2}{\percent} dari hitungan, dan
kedalaman, yang dibaca dari setengah lebar \qty{3.4}{km}, mungkin hingga
\qty{0.1}{km}, masuk dalam bentuk kuadrat, \qty{6}{\percent}; koreksi \qty{3}{\percent}
berada dalam ketidakpastian itu.

\textbf{21.} Pemanasan medan fumarol (Bagian~II) dari waktu ke waktu, peningkatan
gempa bumi kecil di bawah tonjolan, perubahan gas yang
dilepaskan, dan pengangkatan yang diukur dengan penentuan posisi dengan satelit atau sipat datar di
beberapa titik pada lereng.

\textbf{22.} Reservoir di bawah Gunung Halvren membesar sebesar
\textbf{15~juta~\unit[text-series-to-math=true]{m^3}} ($\num{1.51e7}\,\unit{m^3}$).
\end{solution}
